% Generated draft: requires mathematical and source review.
\documentclass[11pt]{article}
\usepackage[margin=25mm]{geometry}
\usepackage{fontspec}
\setmainfont{Latin Modern Roman}
\usepackage{amsmath,amssymb,mathtools,unicode-math}
\setmathfont{Latin Modern Math}
\usepackage{enumitem,needspace,xurl,hyperref}
\setlength{\parindent}{0pt}
\setlength{\parskip}{5pt}
\setlength{\emergencystretch}{3em}
\newcommand{\N}{\mathbb{N}}
\newcommand{\Z}{\mathbb{Z}}
\newcommand{\Q}{\mathbb{Q}}
\newcommand{\R}{\mathbb{R}}
\newcommand{\C}{\mathbb{C}}
\newcommand{\F}{\mathbb{F}}
\newcommand{\A}{\mathbb{A}}
\newcommand{\PP}{\mathbb{P}}
\newcommand{\E}{\mathbb{E}}
\newcommand{\eps}{\varepsilon}
\newcommand{\dd}{\,\mathrm{d}}
\begin{document}
\section*{UnsolvedMath ID: 30000047}
\textbf{Source problem code:} catalogue entry
\subsection{Definitions and mathematical statement}
\textbf{Question.} In electromagnetic inverse scattering, is the factorization-method far-field equation solvable if and only if the sampling point lies in the scatterer \(D\)? The record does not reproduce the far-field equation. To state its intended range test without choosing an unsupported operator, denote the specified factorization operator by \(\mathcal A:H\to Y\), and the far-field sampling datum associated with a point \(z\) and any required polarization \(p\) by \(\phi_{z,p}\in Y\). Once those objects and their hypotheses are fixed, the requested equivalence has the form\[z\in D\quad\Longleftrightarrow\quad\phi_{z,p}\in\operatorname{Ran}(\mathcal A)\quad\Longleftrightarrow\quad\exists g\in H:\ \mathcal A g=\phi_{z,p}.\] The relevant quantifier over \(p\) is not specified: a fixed admissible nonzero polarization, existence of a polarization, and every polarization are distinct formulations. The equation is understood here as exact solvability in the stated spaces, not merely approximate solvability or membership in the closure of the range. The source does not specify how boundary sampling points are treated.\par In a three-dimensional Maxwell setting, far fields are conventionally tangential vector fields on \(S^2\), for example elements of \(L_t^2(S^2)=\{v\in L^2(S^2;\C^3):v(\widehat x)\cdot\widehat x=0\text{ a.e.}\}\), but even this ambient setting is not explicit in the extracted question. The raw far-field operator, a positive operator formed from it, and a fractional power of such an operator must not be interchanged: the precise factorization-method equation determines \(\mathcal A\). The scatterer class, boundary or transmission conditions, frequency assumptions, and any exceptional-spectrum exclusions are also omitted.\par The supplied literature triage checked 2026-08-21 labels the item partially solved. It reports factorization results for important penetrable-medium cases, but says the electromagnetic perfect-conductor analogue matching the intended criterion was still reported open and that no later direct resolution was located in that search. Thus a perfect-conductor focus is suggested by the assessment rather than fixed by the short question itself; results for penetrable media cannot silently be transferred to it. This dated assessment requires independent verification.
\subsection{Short English statement}
The desired test identifies interior sampling points by exact solvability of a specified electromagnetic factorization equation. The equation, polarization quantifiers, and scattering hypotheses are missing. The supplied triage distinguishes established penetrable-medium cases from the intended perfect-conductor analogue.
\subsection{Sources}
\begin{itemize}\item Local UnsolvedMath registry, permanent ID 30000047, code catalogue entry; \emph{Far-Field Solvability in Electromagnetic Inverse Scattering}, Oberwolfach Reports Open Problems.\item \emph{Computational Electromagnetism (2004)}, Oberwolfach Reports, DOI 10.4171/owr/2004/11, pp. 545--632. \url{https://doi.org/10.4171/owr/2004/11}. Proposed year 2004; no proposer named.\item Supplied literature label \emph{Kirsch, Maxwell factorization}: \url{https://doi.org/10.1088/0266-5611/20/6/S08}.\item Supplied literature label \emph{OWR 24/2017}: \url{https://publications.mfo.de/bitstream/handle/mfo/1323/OWP2017_24.pdf}.\item Literature triage checked 2026-08-21, status \texttt{partially\_solved}; record updated 2026-09-16. The \texttt{corrected\_verified} label concerns notation and title normalization, not a supplied operator equation. Verify dated claims independently.\end{itemize}
\end{document}
